% Problem record op_406a00f7c5a9398c. This TeX file is derived from
% database/problems_json/op_406a00f7c5a9398c.json by scripts/sync-tex.mjs; edit the JSON record
% and rerun the script rather than editing this file.
\section{Average-case approximation hardness of squared normalized gaps of random cubic polynomials}
\label{sec:op_406a00f7c5a9398c}

\paragraph{Problem.}
Is it $\#\mathrm{P}$-hard to approximate $\operatorname{ngap}(f)^2$ to
relative multiplicative error $a+o(1)$ on a $b$ fraction of uniformly random
degree-3 polynomials over $\mathbb{F}_2$?
Here $a>0$ and $0<b\leq1$ are constant parameters independent of the number
of variables $n$, specifying the relative error and polynomial fraction.

Write the random polynomial in multilinear form as
\begin{equation}
  f(x)=\sum_{i<j<k}\alpha_{ijk}x_i x_j x_k
       +\sum_{i<j}\beta_{ij}x_i x_j
       +\sum_i\gamma_i x_i \pmod{2},
  \qquad x\in\{0,1\}^n.
  \label{eq:bms-poly-ensemble}
\end{equation}
In Eq.~\eqref{eq:bms-poly-ensemble}, all coefficients are independent uniform
bits. The degree-3 ensemble allows terms of degrees one through three,
without conditioning on a nonzero cubic term. A constant term is omitted
because it changes only the sign of the gap. Define
\begin{equation}
  \operatorname{gap}(f)
    =|\{x\in\{0,1\}^n:f(x)=0\}|-|\{x\in\{0,1\}^n:f(x)=1\}|,
  \qquad
  \operatorname{ngap}(f)=2^{-n}\operatorname{gap}(f).
  \label{eq:bms-poly-gap}
\end{equation}
The target is the square of the normalized gap in Eq.~\eqref{eq:bms-poly-gap}.
An estimate $\widetilde Q_f$ is required to satisfy
\begin{equation}
  \bigl|\widetilde Q_f-\operatorname{ngap}(f)^2\bigr|
    \leq(a+o(1))\operatorname{ngap}(f)^2.
  \label{eq:bms-poly-approximation}
\end{equation}
The $b$ fraction in the question is measured over the coefficient choices
for which Eq.~\eqref{eq:bms-poly-approximation} holds; $o(1)$ tends to zero
as $n\to\infty$.

\subsection*{Status}

Unsolved

\subsection*{Source}

Conjecture 3 of Bremner, Montanaro, and Shepherd, on page 2 of arXiv v2,
states this average-case hardness conjecture with $a=1/4$ and $b=1/24$
\sourcecite{ref:bms-poly-source}{BMS16}.
The present formulation replaces those two numerical constants by the
parameters $a$ and $b$, as requested by the contributor, and retains the
original $o(1)$ term and squared normalized-gap target.
The parameterized formulation is not a verbatim claim of the paper.

\subsection*{Progress}

\begin{itemize}
  \item Appendix D proves worst-case $\#\mathrm{P}$-hardness of approximating
$\operatorname{ngap}(f)^2$ with relative multiplicative error less than $1/2$
\sourcecite{ref:bms-poly-source}{BMS16}.
This worst-case result does not establish hardness on a constant fraction
of uniformly random polynomials.
\end{itemize}

\subsection*{References}

\begin{enumerate}[label={},leftmargin=4.8em,labelsep=0.5em]
  \footnotesize
  \item[\textup{[BMS16]}]\label{ref:bms-poly-source}
  M. J. Bremner, A. Montanaro, and D. J. Shepherd,
"Average-case complexity versus approximate simulation of commuting quantum computations,"
\emph{Physical Review Letters} \textbf{117}, 080501 (2016).
\href{https://doi.org/10.1103/PhysRevLett.117.080501}{doi:10.1103/PhysRevLett.117.080501};
\href{https://arxiv.org/abs/1504.07999v2}{arXiv:1504.07999v2}.
\end{enumerate}

\subsection*{Comment}

The remaining task is to establish average-case approximation hardness
for the uniform polynomial ensemble. The pair $(a,b)$ indexes a family of
questions; hardness for every pair is not asserted. The source's original
parameter choice is $(a,b)=(1/4,1/24)$.
The instance fraction $b$ concerns polynomials, not an algorithm's internal
success probability. The target remains the square of the normalized gap.

\subsection*{Field}

Quantum algorithm

\subsection*{Topic}

Computational complexity; Quantum supremacy; IQP sampling

\subsection*{ID}

\texttt{op\_406a00f7c5a9398c}
